সেট~$A$ সসীম হলে তার আকার একটি স্বাভাবিক সংখ্যা দিয়ে মাপা যায়
এবং ওই সংখ্যাগুলি তুলনা করে আকারের তুলনা করা যায়। সেটগুলি অসীম
হলে ব্যাপারটি আরও জটিল। অসীমতার প্রথম স্তর হলো \emph{অসীম গণনীয়}
সেটগুলির স্তর। সেট~$A$-এর উপাদানগুলিকে একটি \emph{তালিকায়নে},
অর্থাৎ একদিকে অসীম একটি তালিকায়, সাজানো গেলে অশূন্য সেটটি গণনীয়।
অশূন্য সেটের জন্য এর অর্থ একটি সমাপতিত অপেক্ষক~$f\colon \PosInt \to A$
আছে। শূন্য সেটও গণনীয়; তার তালিকায়ন শূন্য তালিকা।
সেটটি গণনীয় অথচ সসীম না হলে সেটি অসীম গণনীয়। কান্টরের
\emph{আঁকাবাঁকা পথের পদ্ধতি} দেখায় যে অসীম গণনীয় সেটগুলির
উপাদানের যুগলগুলির সেটও গণনীয়। এর সাহায্যে দেখানো যায় যে এমনকি
মূলদ সংখ্যার সেট~$\Rat$-ও গণনীয়।

তবে এমন অসীম সেটও আছে যা গণনীয় নয়। এই সেটগুলিকে \emph{অগণনীয়}
বলে। কোনো সেট অগণনীয় দেখানোর দুটি উপায় আছে—সরাসরি একটি
\emph{কর্ণ যুক্তি} দিয়ে, অথবা \emph{হ্রাসকরণ} দিয়ে। কর্ণ যুক্তি
দিতে হলে ধরি যে বিবেচ্য সেট~$A$ গণনীয়, তারপর একটি কাল্পনিক
তালিকায়ন ব্যবহার করে~$A$-এর এমন একটি উপাদান সংজ্ঞায়িত করি, যা
সংজ্ঞা দেওয়ার পদ্ধতির কারণেই তালিকায়নের প্রত্যেক উপাদান থেকে
আলাদা। ফলে তালিকায়নটি গোটা~$A$-এর তালিকায়ন হতে পারে না, এবং
দেখা গেল যে~$A$-এর কোনো তালিকায়নই থাকতে পারে না। হ্রাসকরণে
$A$-এর প্রত্যেক উপাদানকে আগে থেকেই অগণনীয় বলে জানা কোনো সেট~$B$-এর
একটি উপাদানের সঙ্গে সমাপতিতভাবে যুক্ত করে দেখানো হয় যে $A$
অগণনীয়। এমনটি করা গেলে, $A$-এর একটি কাল্পনিক তালিকায়ন
থেকে~$B$-এর একটি তালিকায়ন পাওয়া যেত। যেহেতু $B$ অগণনীয়,
তাই~$A$-এর কোনো তালিকায়ন থাকতে পারে না।

সাধারণভাবে অসীম সেটের আকারের তুলনা করা যায়। $A$ ও $B$-এর মধ্যে
একটি একৈক সমাপতিত অপেক্ষক থাকলে তাদের আকার সমান, অর্থাৎ তারা
\emph{সমসংখ্যক}। $A$ থেকে~$B$-তে একটি একৈক অপেক্ষক থাকলে আমরা
সংজ্ঞা দিতে পারি যে $A$ আকারে~$B$-এর চেয়ে বড় নয়
($\cardle{A}{B}$)। শ্র্যোডার--বার্নস্টাইন উপপাদ্য অনুযায়ী, এতে
বস্তুত অসীম সেটগুলির আকারের একটি ক্রম পাওয়া যায়। সবশেষে,
\emph{কান্টরের উপপাদ্য} বলে, যেকোনো $A$-এর জন্য
$\cardless{A}{\Pow{A}}$। এটি $\Pow{\PosInt}$ অগণনীয় হওয়ার
ফলটির সাধারণীকরণ এবং এটি দেখায় যে অসীমতার স্তর কেবল দুটি নয়,
অসীমসংখ্যক।
